\part{Literature review, novelty and research contribution}

\begin{warnbox}{Read this before Part II}
The repository's literature folder holds three papers. In \code{tracker.md}, the AI-slicing and self-healing literature
reviews are still \emph{open}. The answers below use those three papers plus canonical references; references marked
\textsuperscript{†} are not in the repository and must be checked against the original before citing.
\textbf{Do not cite anything you have not read.} The honest framing for the panel is a targeted review around
three anchors, with a systematic review protocol defined (Q76--77) but not yet fully executed.
\end{warnbox}

\section{Positioning against existing work}

\begin{qa}{59}{What are the three closest research papers to CARS-SF?}
\begin{enumerate}[leftmargin=1.3em]
\item S.~S.~Badawadagi, M.~K.~CM, ``SliceNet: A Hybrid CNN Model for Network Slicing in 5G Communications'', ICEEICT 2024.
\item X.~Limani et~al., ``Demonstrating Isolation Principles for 5G Network Slicing'', IEEE NFV-SDN 2024.
\item A.~M.~Alwakeel, ``Network Slicing and Stand Alone 5G Architectures: Paving the Way for Smart Cities and IoT'', ICEIEC 2025.
\end{enumerate}
Also relevant, cited by [1]: A.~Thantharate et~al., ``DeepSlice'', IEEE UEMCON 2019\textsuperscript{†}.
\end{qa}

\begin{qa}{60}{What problem does each paper solve?}
\begin{itemize}
\item [1] predicts the 3GPP slice type (eMBB/URLLC/mMTC) for a device from a KPI table, using a hybrid CNN, including
for device types not seen in training (reported accuracy \NPaperSnAcc).
\item [2] demonstrates isolation between slices across RAN, transport and 5G core in a real 5G SA proof of concept.
\item [3] surveys 5G SA architecture and NFV-based slicing for smart cities and IoT.
\end{itemize}
\end{qa}

\begin{qa}{61}{What are their assumptions, evaluation environments and limitations?}
\begin{itemize}
\item [1] Tabular KPIs (LTE/5G category, QCI, packet loss, delay; about 31.6k train and 31.6k test rows, 17 columns)
whose label is a deterministic function of declared attributes. There is no network in the loop. Classes are
technology-centric, and evaluation is on a single held-out file.
\item [2] Real hardware; it is a demonstration of isolation mechanisms, not an allocation policy, and has no context
or self-healing.
\item [3] A conceptual and architectural survey with no experimental evaluation of allocation.
\end{itemize}
\end{qa}

\begin{qa}{62}{Which existing approaches already do context-aware network resource allocation?}
Context-aware or intent-based slicing and policy-driven QoS (3GPP PCF policy control, intent-based networking) exist in
general. CARS-SF's specific angle is \emph{socio-temporal} context (school hours, harvest, disaster alerts) selecting
\emph{service} priorities on a rural bottleneck.
\honest{A systematic search for context-aware rural slicing has not been completed. Avoid claiming ``first''.}
\end{qa}

\begin{qa}{63}{Which existing approaches already do fault detection and recovery?}
Fast reroute and link protection (MPLS FRR, BFD-triggered failover\textsuperscript{†}, RFC 5880), OAM (ITU-T Y.1731),
SON self-healing (3GPP SON\textsuperscript{†}) and ETSI ZSM closed-loop automation\textsuperscript{†}. CARS-SF's healing
is conceptually standard protection switching. What is specific is coupling it to slice allocation (backup capacity
ceilings, graceful degradation) and congestion-aware predictive switchover.
\end{qa}

\begin{qa}{64}{Which approaches combine traffic classification with network orchestration?}
ML-based application identification feeding QoS or slicing decisions is common in the literature, and [1] maps
predictions to slices conceptually. CARS-SF closes the loop: a live prediction installs a per-flow HTB rule, and the
SLA outcome is measured (\code{classification/online.py}).
\end{qa}

\begin{qa}{65}{How does CARS-SF differ from a conventional SDN controller with predefined rules?}
An SDN controller is a \emph{mechanism} (it installs match-action rules). CARS-SF is a \emph{policy loop}: it computes
allocations from context and measured demand, enforces them with queueing disciplines (a controller alone does
not shape or schedule), and heals. It could run \emph{on} an SDN controller; the OVS backend uses OpenFlow
\code{set_queue} rules.
\end{qa}

\begin{qa}{66}{How does it differ from static QoS enforcement with Linux tc?}
Static tc is the \texttt{static\_equal} baseline: fixed shares, no context, no demand, no healing. CARS-SF changes HTB
rates at run time (re-planning every 5\,s or on events), adds floors, borrowing ceilings and failure-driven
constraints.
\end{qa}

\begin{qa}{67}{How does it differ from ordinary priority-based scheduling?}
Strict priority starves lower classes and wastes reserved capacity. CARS-SF allocates \emph{weighted} shares capped at
measured demand, redistributes the surplus, and keeps floors. HTB prio only orders borrowing.
The \texttt{proportional} baseline is weighted but demand-blind (Fig.~\ref{fig:alloc}).
\end{qa}

\begin{qa}{68}{How does it differ from existing ML traffic classification?}
\begin{itemize}
\item Labels are \emph{service} classes, not protocols or applications.
\item Features are traffic shape only; QoS outcomes, ports and DSCP are excluded to avoid circularity and shortcuts.
\item Evaluation holds out whole run$\times$UE groups.
\item Predictions are used in the loop.
\end{itemize}
\honest{The traffic is generated by our models, so the F1 is an emulation result (Q334).}
\end{qa}

\begin{qa}{69}{How does it differ from a conventional network-slicing orchestration system?}
Slice orchestrators (lifecycle of slice instances, NFV placement, as in [2] and [3]) create and isolate slices. CARS-SF
assumes the slices exist and decides their \emph{share of a scarce shared link} over time, plus recovery.
These are complementary layers.
\end{qa}

\begin{qa}{70}{Which parts of CARS-SF are established techniques rather than novel?}
HTB, fq\_codel, netem, Mininet, OVS, water-filling allocation, Holt forecasting, random forests, 1D-CNNs, PPO, CEM,
protection switching, Wilcoxon tests. All of them are established.
\end{qa}

\begin{qa}{71}{What exactly is novel about the integration, if the components exist?}
The coupling and the evaluation:
\begin{itemize}
\item healing exposes \emph{constraints} (backup-capacity ceilings, degraded mode) that the allocator honours;
\item the predictor discriminates local congestion from backhaul degradation using the allocator's own queues;
\item the loop is loop-free by construction (quarantine before backup, backup down before primary);
\item outcomes are measured with packet-matched one-way delay on the same platform.
\end{itemize}
\end{qa}

\begin{qa}{72}{Is the novelty algorithmic, architectural, data-centric, methodological or empirical?}
Architectural and empirical first, methodological second (the packet-matched evaluation harness and transfer study).
It is \emph{not} algorithmic: the allocator is a variant of weighted water-filling.
\end{qa}

\begin{qa}{73}{What is the strongest defensible novelty claim?}
``An evaluated closed-loop framework that couples context-driven service slicing with loop-free self-healing on a
shared rural bottleneck. In emulation (L3, including real OVS), it shows that service slicing raises critical-service
SLA compliance (+\NSlaDeltaNoSlicing, $p=\NSlaPNoSlicing$) and that healing cuts backhaul-failure outages to
sub-second, while the choice among reasonable allocation policies and RL adds no significant benefit.''
\end{qa}

\begin{qa}{74}{State that claim without ``intelligent'', ``adaptive'', ``autonomous'' or ``AI-driven''.}
``We built and measured a controller that assigns shares of a rural uplink to four public-service classes according
to time-of-day and emergency events, and that switches to a standby backhaul when the primary fails. Classifying by
service raises clinic traffic's on-time delivery; the switching cuts outages from about 15\,s to under a second in
emulation.''
\end{qa}

\begin{qa}{75}{What is the closest counterexample to your novelty claim?}
A standard deployment that combines operator-configured QoS classes, DiffServ marking and BFD-triggered backup
switchover already gets most of our measured benefit, with no context engine and no ML. Our own results are
consistent with that: slicing presence and fast failover drive the gains. The answer: the contribution is the
integrated, measured evidence for that in the rural setting, plus the coupling (constraints, congestion-aware
prediction), not a superior allocator.
\end{qa}

\begin{qa}{76}{What search methodology did you use for the literature review?}
\honest{Targeted, not systematic: forward and backward snowballing from three anchor papers [1--3].}
Planned protocol to complete before the defense:
\begin{itemize}
\item databases: IEEE Xplore, ACM DL, Scopus;
\item queries: (``network slicing'' OR ``QoS'') AND (rural OR ``developing regions'') AND (context OR ``self-healing'' OR
``traffic classification'');
\item years 2015--2026;
\item PRISMA-style counts of records screened, excluded and included.
\end{itemize}
\end{qa}

\begin{qa}{77}{Which databases, keywords, years and inclusion criteria?}
As in Q76. Inclusion: peer-reviewed, evaluates allocation, healing or classification with measurements or simulation.
Exclusion: vendor white papers, position papers without evaluation (these may be cited as background only).
\end{qa}

\begin{qa}{78}{How do you distinguish peer-reviewed evidence from vendor claims or informal articles?}
Cite only venue-published work for empirical claims, standards (3GPP, IETF RFCs, ITU-T) for definitions, and project
documentation (Open5GS, UERANSIM, Mininet) for tool behaviour, marked as such.
\end{qa}

\begin{qa}{79}{What research gap remains after the closest related work?}
[1] classifies without a network and with technology classes; [2] isolates slices without deciding shares or healing;
[3] has no evaluation. None measures \emph{service}-level outcomes under context changes and backhaul failures on a
shared bottleneck with packet-level evidence.
\end{qa}

\begin{qa}{80}{How would an independent researcher reproduce your comparison?}
\code{make setup data models}, then \code{experiments/run_matrix.py} (\code{configs/experiments/thesis.yaml}), then
\code{aggregate.py} and \code{report.py}. Every run records params, seed and git hash (\code{run_meta.json}).
A sudo-free path (user namespaces) and a Docker image for real OVS exist. Comparing against [1] would need their
dataset; we did not reproduce [1].
\end{qa}

\section{Research contribution and attribution}

\begin{qa}{81}{List your contributions in order of importance.}
\begin{enumerate}[leftmargin=1.3em]
\item The closed-loop framework coupling context-driven service slicing and loop-free self-healing (with event-driven
failover and congestion-aware prediction).
\item The evaluation methodology: endpoint packet matching, a 340-run matrix, group-held-out ML evaluation.
\item The empirical findings: slicing presence matters, allocation-policy choice and RL do not, healing works,
synthetic data does not transfer.
\item SliceNet on measured traffic and its in-loop use.
\item A 5G-core (Open5GS) enforcement demonstration.
\end{enumerate}
\end{qa}

\begin{qa}{82}{Which contribution can be evaluated independently?}
The measurement methodology (validated by exact propagation floors \NFloorEdge/\NFloorCore\,ms and unit tests on
synthetic pcaps), SliceNet (offline F1), and the transfer study.
\end{qa}

\begin{qa}{83}{Which contribution depends on another being correct?}
All network-outcome claims depend on the measurement pipeline. In-loop SliceNet depends on SliceNet and the QoS
enforcer. RL results depend on simulator calibration.
\end{qa}

\begin{qa}{84}{What does RSNSD contribute beyond a collection of synthetic records?}
A service-semantic schema (89 columns, provenance-tagged), scenario-aware generation (19 scenarios, 32 applications)
and a reproducible generator.
\honest{Its value as \emph{training data for measured traffic} is refuted by our own transfer study (zero-shot F1
\NTrZero). Its defensible role is as a schema and bootstrap; the measured dataset supersedes it.}
\end{qa}

\begin{qa}{85}{What does SliceNet contribute beyond training a classifier on RSNSD?}
It is trained on \emph{measured} windows (\NSnRows{} from \NSnRunsCount{} runs), uses only shape features at a single
vantage point, is evaluated on held-out run$\times$UE groups (RF \NSnRfFone, CNN \NSnCnnFone), and is deployed in the
loop for a mixed-service UE.
\end{qa}

\begin{qa}{86}{What does the allocator contribute beyond a predefined priority order?}
Demand-awareness (measured offered load, elastic probe-up, Holt forecast), floors, surplus redistribution, and
healing constraints.
\honest{The main study does not show a significant utility gain over simpler policies; the gain appears only when
HTB borrowing is disabled (ablation) or at $\ge$1$\times$ load.}
\end{qa}

\begin{qa}{87}{What does self-healing contribute beyond executing a recovery action?}
\begin{itemize}
\item detection in milliseconds through an event-driven link watcher;
\item diagnosis (link vs core switch, by elimination);
\item loop-free ordering;
\item hold-down plus OAM-gated fail-back;
\item suppression of its own link events;
\item prediction from RTT trends, with congestion discrimination;
\item allocation constraints during the failure.
\end{itemize}
\end{qa}

\begin{qa}{88}{What does the closed loop contribute beyond chaining modules?}
Shared state and feedback. The allocator sees healing state; the predictor sees allocator queues; the demand estimator
sees enforcement saturation (probe-up); re-planning is event-triggered. Two real bugs were found \emph{only} at the
seams: bufferbloat on the backup link, and false failover under AQM-less queues.
\end{qa}

\begin{qa}{89}{Which contribution has the strongest experimental evidence?}
Self-healing (deterministic, large effect: outage \NOutCutFullMean{} vs \NOutCutNoneMean\,s, replicated on real OVS),
and slicing vs no slicing for healthcare SLA (Bonferroni-significant).
\end{qa}

\begin{qa}{90}{Which contribution currently has the weakest evidence?}
Context awareness (no measurable ablation effect), RL (null result), and anything about real rural
deployment (no L4 evidence).
\end{qa}

\begin{qa}{91}{Which parts did you personally design, implement, evaluate or validate?}
\textit{Answer truthfully for your own case.} The repository history attributes the framework, registries, traffic
models, measurement pipeline, controller, healing engine, SliceNet pipeline, experiment tooling and 5G deployment to
this project. Be ready to state which were AI-assisted and how you verified them.
\end{qa}

\begin{qa}{92}{Which parts were adapted from existing tools or libraries?}
Mininet, Open vSwitch, Linux tc (HTB, fq\_codel, netem), tcpdump and dpkt, scikit-learn, PyTorch, stable-baselines3 and
Gymnasium, Open5GS and UERANSIM (Docker images by Gradiant), pandas and SciPy.
\end{qa}

\begin{qa}{93}{What would remain as a contribution if the ML component were removed?}
Almost everything that is significant: service slicing with healing, the measurement methodology, and the
empirical findings. ML contributes the transfer result and the mixed-UE capability only.
\end{qa}

\begin{qa}{94}{What would remain if the recovery mechanism were replaced by a static policy?}
A static ``always use the backup when the primary is down'' rule \emph{is} our reactive core. What remains specific is
event-driven speed, safe fail-back, fault localisation, predictive switchover (outage \NOutDegFullMean{} vs
\NOutDegReactiveMean\,s reactive-only) and the allocation coupling.
\end{qa}

\begin{qa}{95}{What would remain if RSNSD were replaced with a different dataset?}
Everything except the transfer study. The in-loop classifier already uses the measured dataset, not RSNSD.
\end{qa}

\begin{qa}{96}{What would persuade a skeptical reviewer that this is more than software integration?}
Falsifiable pre-registered hypotheses with reported negative results; packet-level measurement validated against
ground truth (exact delay floors); statistics (CIs, paired tests, Bonferroni); integration bugs found and fixed
(busy-wait timing distortion, partial-window bias, uplink-only utilisation, false failover) with all runs redone.
That is research method, not plumbing.
\end{qa}
